\section{Edge-Middle Reduction}\label{sec:peps_edge_middle}

Fixing an edge $e=(u,v)$ and blocking the whole complementary middle region
$V\setminus\{u,v\}$ presents the state coefficient as a three-site contraction:
the endpoint tensors at $u$ and $v$ joined through one blocked middle tensor.
Inserting the identity on the bond $e$ recovers the original coefficient, so a
matrix inserted on $e$ records the difference between two tensors that generate
the same state. Injectivity of this middle tensor provides the foundation
used by the edge-kernel argument of Section~\ref{sec:peps_edge_kernel_gauges}.

\begin{definition}[Open middle tensor at an edge]%
    \label{def:peps_edgeOpenMiddleWeight}
    \lean{TNLean.PEPS.edgeOpenMiddleWeight}
    \leanok
    \uses{def:peps_edgeInsertedBoundaryConfig, def:peps_edgeMiddleVertices}
    For an edge $e=(u,v)$, $M_e=V\setminus\{u,v\}$, and residual endpoint
    configurations $(\lambda,\rho)$, the open middle tensor is
    $C^{\mathrm{open}}_{A,M_e}(\lambda,\rho;\sigma)
        =\sum_{\eta\in\mathcal C_e(\lambda,\rho)}\prod_{x\in M_e}A_x(\eta|_x,\sigma_x)$,
    summed over virtual configurations $\eta$ on $E\setminus\{e\}$ matching
    $(\lambda,\rho)$.
\end{definition}

\begin{definition}[Middle physical configurations at an edge]%
    \label{def:peps_edgeMiddlePhysicalConfig}
    \lean{TNLean.PEPS.EdgeMiddlePhysicalConfig}
    \lean{TNLean.PEPS.edgeMiddlePhysicalConfigOf}
    \leanok
    \uses{def:peps_edgeMiddleVertices}
    For $e=(u,v)$, a middle physical configuration is an assignment
    $\tau:M_e\to\{0,\ldots,d-1\}$, and a global configuration $\sigma$
    restricts to $\sigma|_{M_e}$.
\end{definition}

\begin{definition}[Middle tensor with middle physical index]%
    \label{def:peps_edgeMiddleWeightOn}
    \lean{TNLean.PEPS.edgeMiddleWeightOn}
    \lean{TNLean.PEPS.edgeOpenMiddleWeightOn}
    \leanok
    \uses{def:peps_edgeMiddlePhysicalConfig, def:peps_edgeMiddleConfig, def:peps_edgeOpenMiddleWeight}
    For the middle tensor in the three-site chain of
    \cite[Section~3]{Molnar2018NormalPEPS}, write
    $C_{A,M_e}(\beta,\tau)$ and $C^{\mathrm{open}}_{A,M_e}(\lambda,\rho;\tau)$
    for the ordinary and open middle contractions evaluated on a middle
    physical configuration $\tau:M_e\to\{0,\ldots,d-1\}$.
\end{definition}

\begin{theorem}[Middle restriction of the blocked tensor]%
    \label{thm:peps_edgeMiddleWeight_eq_on}
    \lean{TNLean.PEPS.edgeMiddleWeight_eq_on}
    \lean{TNLean.PEPS.edgeOpenMiddleWeight_eq_on}
    \leanok
    \uses{def:peps_edgeMiddleWeightOn, def:peps_edgeMiddlePhysicalConfig}
    For every global physical configuration $\sigma$,
    $C_{A,M_e}(\beta,\sigma)=C_{A,M_e}(\beta,\sigma|_{M_e})$ and
    $C^{\mathrm{open}}_{A,M_e}(\lambda,\rho;\sigma)=C^{\mathrm{open}}_{A,M_e}(\lambda,\rho;\sigma|_{M_e})$.
\end{theorem}
\begin{proof}\leanok
    Each vertex product ranges over $x\in M_e$, where $\sigma_x=(\sigma|_{M_e})_x$.
\end{proof}

\begin{definition}[Restricting middle configurations away from the edge]%
    \label{def:peps_edgeMiddleConfigToOpenMiddleConfig}
    \lean{TNLean.PEPS.edgeMiddleConfigToOpenMiddleConfig}
    \leanok
    \uses{def:peps_edgeMiddleConfig, def:peps_edgeOpenMiddleWeight}
    An ordinary middle configuration $\eta$ compatible with
    $\beta=(a,\lambda,\rho)$ restricts to
    $\eta|_{E\setminus\{e\}}\in\mathcal C_e(\lambda,\rho)$.
\end{definition}

\begin{definition}[Restoring the distinguished edge index]%
    \label{def:peps_edgeOpenMiddleConfigToMiddleConfig}
    \lean{TNLean.PEPS.edgeOpenMiddleConfigToMiddleConfig}
    \leanok
    \uses{def:peps_edgeBoundaryConfig, def:peps_edgeMiddleConfigToOpenMiddleConfig}
    Conversely, extending $\eta\in\mathcal C_e(\lambda,\rho)$ by $\eta(e)=a$
    yields a middle configuration compatible with $\beta=(a,\lambda,\rho)$.
\end{definition}

\begin{definition}[Open-middle reindexing equivalence]%
    \label{def:peps_edgeMiddleConfigEquivOpenMiddleConfig}
    \lean{TNLean.PEPS.edgeMiddleConfigEquivOpenMiddleConfig}
    \leanok
    \uses{def:peps_edgeMiddleConfigToOpenMiddleConfig, def:peps_edgeOpenMiddleConfigToMiddleConfig}
    For fixed $\beta=(a,\lambda,\rho)$, the maps of
    Definitions~\ref{def:peps_edgeMiddleConfigToOpenMiddleConfig}
    and~\ref{def:peps_edgeOpenMiddleConfigToMiddleConfig} form a bijection
    $\{\eta:\partial_e(\eta)=\beta\}\simeq\mathcal C_e(\lambda,\rho)$.
\end{definition}

\begin{theorem}[Middle-indexed open tensor reindexing]%
    \label{thm:peps_edgeMiddleWeightOn_eq_edgeOpenMiddleWeightOn}
    \lean{TNLean.PEPS.edgeMiddleWeightOn_eq_edgeOpenMiddleWeightOn}
    \leanok
    \uses{def:peps_edgeMiddleWeightOn, def:peps_edgeMiddleConfigEquivOpenMiddleConfig}
    For $\beta=(a,\lambda,\rho)$ and $\tau:M_e\to\{0,\ldots,d-1\}$,
    $C_{A,M_e}((a,\lambda,\rho),\tau)=C^{\mathrm{open}}_{A,M_e}(\lambda,\rho;\tau)$.
\end{theorem}
\begin{proof}\leanok
    Reindex the sum by the bijection of
    Definition~\ref{def:peps_edgeMiddleConfigEquivOpenMiddleConfig}; since
    $e\not\ni x$ for $x\in M_e$, the local factors $A_x(\eta|_x,\tau_x)$ agree.
\end{proof}

\begin{definition}[Edge insertion coefficient]%
    \label{def:peps_edgeInsertedCoeff}
    \lean{TNLean.PEPS.edgeInsertedCoeff}
    \leanok
    \uses{def:peps_edgeInsertedBoundaryConfig, def:peps_edgeOpenMiddleWeight}
    For an edge $e=(u,v)$ and $X\in\MN{D_e}$, the edge insertion coefficient is
    % Formula: C_A(e,X;sigma) is the sum below over inserted boundary data beta.
    % Ink-to-index: there is no retained ink.  The fixed e, sigma and X coefficient is
    % a scalar, with boundary signature (V,P)=(0,0); the old three-open-leg picture
    % instead drew the unfixed family.  No exact standalone picture occurs in
    % Papers/1804.04964/paper_normal.tex; see the Section 3 route recorded in
    % docs/paper-gaps/peps_injective_ft_section3_route.tex.
    \begin{align}
        C_A(e,X;\sigma)
         & =
        \sum_{\beta}
        A_u(\beta_u,\sigma_u)
        X_{\beta_{e,u},\beta_{e,v}}
        C^{\mathrm{open}}_{A,M_e}
        (\beta_{\widehat e,u},\beta_{\widehat e,v};
        \sigma|_{M_e})
        A_v(\beta_v,\sigma_v).
        \notag
    \end{align}
\end{definition}

\begin{theorem}[Middle tensor reindexing for an opened edge]%
    \label{thm:peps_edgeMiddleWeight_eq_edgeOpenMiddleWeight}
    \lean{TNLean.PEPS.edgeMiddleWeight_eq_edgeOpenMiddleWeight}
    \leanok
    \uses{def:peps_edgeOpenMiddleWeight, def:peps_edgeMiddleConfigEquivOpenMiddleConfig}
    For $\beta=(a,\lambda,\rho)$,
    $C_{A,M_e}(\beta,\sigma)=C^{\mathrm{open}}_{A,M_e}(\lambda,\rho;\sigma|_{M_e})$.
\end{theorem}
\begin{proof}\leanok
    Reindex by Definition~\ref{def:peps_edgeMiddleConfigEquivOpenMiddleConfig}
    and compare the local factors $A_x(\eta|_x,\sigma_x)$ for $x\in M_e$.
\end{proof}

\begin{theorem}[Diagonal summand of the identity insertion]%
    \label{thm:peps_edgeInsertedCoeff_identity_diagonal_summand}
    \lean{TNLean.PEPS.edgeInsertedCoeff_identity_diagonal_summand}
    \leanok
    \uses{def:peps_edgeBoundaryToInsertedBoundaryConfig, thm:peps_edgeMiddleWeight_eq_edgeOpenMiddleWeight}
    On the diagonal datum $\Delta(\beta)=(a,a,\lambda,\rho)$ of
    Definition~\ref{def:peps_edgeBoundaryToInsertedBoundaryConfig}, the
    $\Id$-inserted summand equals the ordinary edge-blocked summand at $\beta$.
\end{theorem}
\begin{proof}\leanok
    Substitute $\Id_{a,a}=1$ and
    $C_{A,M_e}(\beta,\sigma)=C^{\mathrm{open}}_{A,M_e}(\lambda,\rho;\sigma|_{M_e})$
    from Theorem~\ref{thm:peps_edgeMiddleWeight_eq_edgeOpenMiddleWeight}.
\end{proof}

\begin{theorem}[Off-diagonal summands of the identity insertion]%
    \label{thm:peps_edgeInsertedCoeff_identity_offDiagonal_summand}
    \lean{TNLean.PEPS.edgeInsertedCoeff_identity_offDiagonal_summand}
    \leanok
    \uses{def:peps_edgeInsertedCoeff}
    For $\beta=(a,b,\lambda,\rho)$ with $a\ne b$, the $\Id$-inserted summand in
    Definition~\ref{def:peps_edgeInsertedCoeff} vanishes.
\end{theorem}
\begin{proof}\leanok
    For $a\ne b$, $\Id_{a,b}=0$.
\end{proof}

\begin{theorem}[Identity insertion recovers the edge-blocked coefficient]%
    \label{thm:peps_edgeInsertedCoeff_identity}
    \lean{TNLean.PEPS.edgeInsertedCoeff_identity}
    \leanok
    \uses{def:peps_edgeInsertedCoeff, def:peps_edgeBlockedCoeff, def:peps_edgeBoundaryToInsertedBoundaryConfig, def:peps_edgeBoundaryConfigEquivDiagonalInsertedBoundaryConfig, thm:peps_edgeInsertedCoeff_identity_diagonal_summand, thm:peps_edgeInsertedCoeff_identity_offDiagonal_summand}
    For every $e$ and $\sigma$, $C_A(e,\Id;\sigma)=C_A(e;\sigma)$.
\end{theorem}
\begin{proof}\leanok
    Splitting the sum in Definition~\ref{def:peps_edgeInsertedCoeff} into
    $a=b$ and $a\ne b$ gives
    \begin{align}
        C_A(e,\Id;\sigma)
         & =
        \sum_{\beta:\,a=b}
        A_u(\beta_u,\sigma_u)
        C^{\mathrm{open}}_{A,M_e}(\beta,\sigma|_{M_e})
        A_v(\beta_v,\sigma_v)
        +\sum_{\beta:\,a\ne b}0,
        \notag
    \end{align}
    which reindexes to $C_A(e;\sigma)$ via
    Definition~\ref{def:peps_edgeBoundaryConfigEquivDiagonalInsertedBoundaryConfig}
    and Theorems~\ref{thm:peps_edgeInsertedCoeff_identity_diagonal_summand}
    and~\ref{thm:peps_edgeInsertedCoeff_identity_offDiagonal_summand}.
\end{proof}

\begin{definition}[Blocked middle tensor at an edge]%
    \label{def:peps_edgeMiddleWeight}
    \lean{TNLean.PEPS.edgeMiddleWeight}
    \leanok
    \uses{def:peps_edgeBoundaryMatches, def:peps_edgeMiddleConfig, def:peps_edgeMiddleVertices}
    For boundary data $\beta$, the blocked middle tensor is
    $C_{A,M_e}(\beta,\sigma)=\sum_{\eta:\partial_e(\eta)=\beta}\prod_{x\in M_e}A_x(\eta|_x,\sigma_x)$,
    presenting the contraction as a three-site chain:
    % Formula (paper_normal.tex:981--1009): choose edge (1,5), set
    % A'_1=A_5, A'_2=A_1, and block vertices {2,3,4} into A'_3; the resulting
    % cyclic word is A'_1 A'_2 A'_3 in that order.
    % Ink-to-index: the left panel has the source five-cycle and chord (2,5),
    % with singleton endpoint regions A'_1,A'_2 and middle region A'_3.  The
    % right panel contracts the three induced virtual bonds cyclically.
    % Contraction and boundary signature: the graph panel has all six virtual
    % edges contracted and five physical legs, (V,P)=(0,5); blocking groups
    % those physical factors into three legs on the cyclic word, (0,3).
    % Source verdict: exact source blocking and cyclic order.  The legacy
    % catalogue word A'_1,A'_3,A'_2 did not match the source.
    \begin{center}
        \tenkzkernel
        % The five vertices sit at the stations of a circle frame, the
        % cycle as the ring and the chord as the one extra edge; the three
        % blocked regions are contours over the station selections.
        \begin{tenkz}[rows={wire}, cols=5, frame=circle, bonds=none]
            \tn[at=(1,1), name=blockAone, label pos=145,
                ports={0:virtual, 180:virtual, 90:physical}]{A_1}
            \tn[at=(1,2), name=blockAtwo, label pos=145,
                ports={0:virtual, 180:virtual, -72:virtual, 90:physical}]{A_2}
            \tn[at=(1,3), name=blockAthree, label pos=145,
                ports={0:virtual, 180:virtual, 90:physical}]{A_3}
            \tn[at=(1,4), name=blockAfour, label pos=145,
                ports={0:virtual, 180:virtual, 90:physical}]{A_4}
            \tn[at=(1,5), name=blockAfive, label pos=145,
                ports={0:virtual, 180:virtual, 252:virtual, 90:physical}]{A_5}
            \tnwire{blockAone.0}{blockAtwo.180}
            \tnwire[name=blockTwoThree]{blockAtwo.0}{blockAthree.180}
            \tnwire[name=blockThreeFour]{blockAthree.0}{blockAfour.180}
            \tnwire{blockAfour.0}{blockAfive.180}
            \tnwire{blockAfive.0}{blockAone.180}
            \tnwire{blockAtwo.-72}{blockAfive.252}
            \tnmark[form=enclosure, species=selected, label pos=90]{blockAone}{$A'_2$}
            \tnmark[form=enclosure, species=secondary, label pos=90]{blockAfive}{$A'_1$}
            % The hull of the three stations contains the two wires
            % between them, so the contour needs only the atoms.
            % The bearing is read in the tilted hull's own axes; its north
            % faces the page's lower-right corner here.
            \tnmark[form=enclosure, species=collar, label pos=n]
            {{blockAtwo, blockAthree, blockAfour}}{$A'_3$}
        \end{tenkz}
        $\quad\rightsquigarrow\quad$
        \begin{tenkz}[rows={ket}, cols=3, physical=up, boundary=periodic]
            \tn[label pos=s]{A'_1} & \tn[label pos=s]{A'_2} &
            \tn[label pos=s]{A'_3}
        \end{tenkz}
    \end{center}
\end{definition}

\begin{definition}[Edge-blocked PEPS coefficient]%
    \label{def:peps_edgeBlockedCoeff}
    \lean{TNLean.PEPS.edgeBlockedCoeff}
    \leanok
    \uses{def:peps_edgeLeftLocalConfig, def:peps_edgeMiddleWeight, def:peps_edgeRightLocalConfig}
    At $e=(u,v)$, the edge-blocked coefficient is
    $C_A(e;\sigma)=\sum_{\beta}A_u(\beta_u,\sigma_u)C_{A,M_e}(\beta,\sigma)A_v(\beta_v,\sigma_v)$.
\end{definition}

\begin{theorem}[Edge-blocked coefficient equals the PEPS coefficient]%
    \label{thm:peps_edgeBlockedCoeff_eq_stateCoeff}
    \lean{TNLean.PEPS.edgeBlockedCoeff_eq_stateCoeff}
    \leanok
    \uses{def:peps_edgeBlockedCoeff, thm:peps_stateCoeff_splitAtEdge, def:peps_virtualConfigEquivEdgeBoundary}
    For every edge $e$, $C_A(e;\sigma)=\Coeff_A(\sigma)$.
\end{theorem}
\begin{proof}\leanok
    \uses{def:peps_virtualConfigEquivEdgeBoundary, thm:peps_edgeLeftLocalConfig_eq_of_boundaryMatches, thm:peps_edgeRightLocalConfig_eq_of_boundaryMatches}
    Reindexing Theorem~\ref{thm:peps_stateCoeff_splitAtEdge} by
    $\omega\leftrightarrow(\beta,\omega_{M_e})$ from
    Definition~\ref{def:peps_virtualConfigEquivEdgeBoundary} and using
    $\omega|_u=\beta_u$, $\omega|_v=\beta_v$ from
    Theorems~\ref{thm:peps_edgeLeftLocalConfig_eq_of_boundaryMatches}
    and~\ref{thm:peps_edgeRightLocalConfig_eq_of_boundaryMatches} gives
    \begin{align}
        \sum_{\omega}
        A_u(\omega|_u,\sigma_u)
        \left(\prod_{x\in M_e}A_x(\omega|_x,\sigma_x)\right)
        A_v(\omega|_v,\sigma_v)
         & =
        \sum_{\beta}
        A_u(\beta_u,\sigma_u)
        C_{A,M_e}(\beta,\sigma|_{M_e})
        A_v(\beta_v,\sigma_v).
        \notag
    \end{align}
\end{proof}

\begin{theorem}[Identity insertion equals the PEPS coefficient]%
    \label{thm:peps_edgeInsertedCoeff_identity_eq_stateCoeff}
    \lean{TNLean.PEPS.edgeInsertedCoeff_identity_eq_stateCoeff}
    \leanok
    \uses{thm:peps_edgeInsertedCoeff_identity, thm:peps_edgeBlockedCoeff_eq_stateCoeff}
    For every $e$ and $\sigma$, $C_A(e,\Id;\sigma)=\Coeff_A(\sigma)$.
\end{theorem}
\begin{proof}\leanok
    Combine Theorems~\ref{thm:peps_edgeInsertedCoeff_identity}
    and~\ref{thm:peps_edgeBlockedCoeff_eq_stateCoeff}:
    $C_A(e,\Id;\sigma)=C_A(e;\sigma)=\Coeff_A(\sigma)$.
\end{proof}

\begin{theorem}[Same state gives the same edge-blocked coefficients]%
    \label{thm:peps_sameState_edgeBlockedCoeff_eq}
    \lean{TNLean.PEPS.SameState.edgeBlockedCoeff_eq}
    \leanok
    \uses{def:peps_same_state, thm:peps_edgeBlockedCoeff_eq_stateCoeff}
    If $\Coeff_A=\Coeff_B$, then $C_A(e;\sigma) = C_B(e;\sigma)$ for all $e,\sigma$.
\end{theorem}
\begin{proof}\leanok
    By Theorem~\ref{thm:peps_edgeBlockedCoeff_eq_stateCoeff},
    $C_A(e;\sigma)=\Coeff_A(\sigma)=\Coeff_B(\sigma)=C_B(e;\sigma)$.
\end{proof}

\begin{theorem}[Same state gives the same identity insertions]%
    \label{thm:peps_sameState_edgeInsertedCoeff_identity_eq}
    \lean{TNLean.PEPS.SameState.edgeInsertedCoeff_identity_eq}
    \leanok
    \uses{def:peps_same_state, thm:peps_edgeInsertedCoeff_identity_eq_stateCoeff}
    If $\Coeff_A=\Coeff_B$, then $C_A(e,\Id;\sigma) = C_B(e,\Id;\sigma)$ for all
    $e,\sigma$.
\end{theorem}
\begin{proof}\leanok
    By Theorem~\ref{thm:peps_edgeInsertedCoeff_identity_eq_stateCoeff},
    $C_A(e,\Id;\sigma)=\Coeff_A(\sigma)=\Coeff_B(\sigma)=C_B(e,\Id;\sigma)$.
\end{proof}

\begin{definition}[Middle tensor family at an edge]%
    \label{def:peps_edgeMiddleTensorFamily}
    \lean{TNLean.PEPS.edgeMiddleTensorFamily}
    \lean{TNLean.PEPS.EdgeMiddleTensorInjective}
    \leanok
    \uses{def:peps_edgeMiddleWeightOn}
    At $e=(u,v)$ with $M_e=V\setminus\{u,v\}$, the edge-middle tensor family is
    $T_{A,M_e}=(T_{A,M_e}^{\rho})_{\rho=(\beta_{\mathrm{L}},\beta_{\mathrm{R}})}$,
    where $T_{A,M_e}^{\rho}(\tau)=C_{\mathrm{mid}}(\tau;\beta_{\mathrm{L}},\beta_{\mathrm{R}})$,
    and its injectivity is $\inj(T_{A,M_e})$.
\end{definition}

\begin{lemma}[Middle tensor vanishes for an empty complement configuration]%
    \label{lem:peps_edgeMiddleTensorFamily_eq_zero_of_isEmpty}
    \lean{TNLean.PEPS.edgeMiddleTensorFamily_eq_zero_of_isEmpty}
    \leanok
    \uses{def:peps_edgeMiddleTensorFamily}
    If the complement virtual configuration space $\mc C_e$ at $e$ is empty,
    then $T_{A,M_e}^{\rho}=0$ for every $\rho$.
\end{lemma}
\begin{proof}\leanok
    Each entry of $T_{A,M_e}^{\rho}$ in Definition~\ref{def:peps_edgeMiddleTensorFamily}
    is a sum over $\mc C_e=\varnothing$, hence $0$.
\end{proof}

\begin{lemma}[Failure of middle injectivity for an empty complement configuration]%
    \label{lem:peps_not_edgeMiddleTensorInjective_of_isEmpty}
    \lean{TNLean.PEPS.not_edgeMiddleTensorInjective_of_isEmpty}
    \leanok
    \uses{def:peps_edgeMiddleTensorFamily, lem:peps_edgeMiddleTensorFamily_eq_zero_of_isEmpty}
    If $\mc C_e=\varnothing$ and the residual boundary-label set is nonempty,
    then $\neg\inj(T_{A,M_e})$.
\end{lemma}
\begin{proof}\leanok
    Choose a boundary label $\rho$; by
    Lemma~\ref{lem:peps_edgeMiddleTensorFamily_eq_zero_of_isEmpty},
    $T_{A,M_e}^{\rho}=0$, which precludes linear independence.
\end{proof}

\begin{definition}[Edge-blocked three-site injectivity]%
    \label{def:peps_edgeBlockedThreeSiteInjective}
    \lean{TNLean.PEPS.EdgeBlockedThreeSiteInjective}
    \leanok
    \uses{def:peps_localTensorMap, def:peps_edgeMiddleTensorFamily}
    At $e=(u,v)$ with $M_e=V\setminus\{u,v\}$, edge-blocked three-site
    injectivity is the conjunction
    $\inj(T_{A,u})\wedge\inj(T_{A,M_e})\wedge\inj(T_{A,v})$.
\end{definition}

\begin{definition}[Singleton blocked region]%
    \label{def:peps_singletonRegionTensor}
    \lean{TNLean.PEPS.SingletonRegionPhysicalConfig,
        TNLean.PEPS.singletonRegionTensorFamily,
        TNLean.PEPS.SingletonRegionTensorInjective,
        TNLean.PEPS.SingletonRegionInjectivityComparison}
    \leanok
    \uses{def:peps_regionInjectivityData}
    For a singleton block $\{v\}$, the blocked tensor is
    $T_{A,\{v\}}(\eta,\sigma)=A_v(\eta,\sigma)$ with injectivity predicate
    $\inj(T_{A,\{v\}})$, and singleton comparison for a region-injectivity
    predicate $\kappa$ is $\forall v,\ \inj(T_{A,\{v\}})\Longrightarrow\kappa(\{v\})$.
\end{definition}

\begin{theorem}[Singleton comparison for a vertex-injective tensor]%
    \label{thm:peps_singletonRegionInjectivityComparison_of_isVertexInjective}
    \lean{TNLean.PEPS.singletonRegionInjectivityComparison_of_isVertexInjective}
    \leanok
    \uses{def:IsVertexInjective, def:peps_singletonRegionTensor, def:peps_regionInjectivityDataOf, thm:peps_regionBlockedTensorInjective_of_isVertexInjective}
    If $A$ is vertex-injective and every bond dimension is positive, then the
    concrete predicate $\kappa(R)=\inj(T_{A,R})$ of
    Definition~\ref{def:peps_regionInjectivityDataOf} satisfies
    $\forall v,\ \inj(T_{A,\{v\}})\Longrightarrow\kappa(\{v\})$.
\end{theorem}
\begin{proof}\leanok
    By Theorem~\ref{thm:peps_regionBlockedTensorInjective_of_isVertexInjective},
    $\kappa(\{v\})$ holds unconditionally for every $v\in V$.
\end{proof}

\begin{theorem}[Vertex injectivity gives singleton-block injectivity]%
    \label{thm:peps_singletonRegionTensorInjective}
    \lean{TNLean.PEPS.IsVertexInjective.singletonRegionTensorInjective}
    \leanok
    \uses{def:IsVertexInjective, def:peps_singletonRegionTensor}
    If $A$ is vertex-injective, then $\inj(T_{A,\{v\}})$ holds for every $v\in V$.
\end{theorem}
\begin{proof}\leanok
    By Definition~\ref{def:peps_singletonRegionTensor},
    $T_{A,\{v\}}=A_v$, so $\inj(T_{A,\{v\}})\iff\inj(A_v)$, which holds by
    Definition~\ref{def:IsVertexInjective}.
\end{proof}

\begin{theorem}[Singleton comparison gives the region base case]%
    \label{thm:peps_singletonRegionInjectivityFromVertexInjectivity}
    \lean{TNLean.PEPS.SingletonRegionInjectivityFromVertexInjectivity,
        TNLean.PEPS.SingletonRegionInjectivityFromVertexInjectivity.ofSingletonComparison}
    \leanok
    \uses{def:IsVertexInjective, def:peps_singletonRegionTensor, thm:peps_singletonRegionTensorInjective}
    If $\forall v,\ \inj(T_{A,\{v\}})\Longrightarrow\kappa(\{v\})$, then
    vertex injectivity of $A$ implies $\forall v,\ \kappa(\{v\})$.
\end{theorem}
\begin{proof}\leanok
    Compose Theorem~\ref{thm:peps_singletonRegionTensorInjective} with the
    hypothesis $\inj(T_{A,\{v\}})\Longrightarrow\kappa(\{v\})$.
\end{proof}

\begin{theorem}[Finite regions from singleton regions and union closure]%
    \label{thm:peps_finiteRegionInjectivityFromSingletons}
    \lean{TNLean.PEPS.RegionInjectivityUnionClosure.finset_injective_of_singletons}
    \leanok
    \uses{thm:peps_singletonRegionInjectivityFromVertexInjectivity, lem:peps_regionInjectivityUnionClosure_finite}
    If $\kappa$ satisfies $\forall v,\ \inj(T_{A,\{v\}})\Longrightarrow\kappa(\{v\})$
    and $\kappa(R)\wedge\kappa(S)\Longrightarrow\kappa(R\cup S)$, and $A$ is
    vertex-injective, then $\kappa(R)$ holds for every finite $R\ne\varnothing$.
\end{theorem}
\begin{proof}\leanok
    Writing $R=\bigcup_{v\in R}\{v\}$,
    Theorem~\ref{thm:peps_singletonRegionInjectivityFromVertexInjectivity} gives
    $\kappa(\{v\})$ for all $v\in R$, and
    Lemma~\ref{lem:peps_regionInjectivityUnionClosure_finite} yields $\kappa(R)$.
\end{proof}

\begin{definition}[Region injectivity gives the edge-middle tensor]%
    \label{def:peps_edgeMiddleRegionInjectivityComparison}
    \lean{TNLean.PEPS.EdgeMiddleRegionInjectivityComparison}
    \leanok
    \uses{def:peps_regionInjectivityData, def:peps_edgeMiddleTensorFamily}
    To isolate the finite-region contraction claim in
    \cite[Section~3]{Molnar2018NormalPEPS}, edge-middle region comparison for
    $\kappa$ at $e=(u,v)$ with $M_e=V\setminus\{u,v\}$ is the implication
    $\kappa(M_e)\Longrightarrow\inj(T_{A,M_e})$.
\end{definition}

\begin{theorem}[Endpoint injectivity in the edge-blocked chain]%
    \label{thm:peps_edgeBlockedEndpointTensorMaps_injective}
    \lean{TNLean.PEPS.IsVertexInjective.edgeBlockedEndpointTensorMaps_injective}
    \leanok
    \uses{def:IsVertexInjective}
    If $A$ is vertex-injective, then for every edge $e=(u,v)$,
    $\inj(T_{A,u})\wedge\inj(T_{A,v})$.
\end{theorem}
\begin{proof}\leanok
    Specialize Definition~\ref{def:IsVertexInjective} to the endpoints $u$ and $v$.
\end{proof}

\begin{theorem}[Reducing edge-blocked injectivity to the middle block]%
    \label{thm:peps_edgeBlockedThreeSiteInjective_of_middle}
    \lean{TNLean.PEPS.IsVertexInjective.edgeBlockedThreeSiteInjective_of_middle}
    \leanok
    \uses{def:IsVertexInjective, def:peps_edgeBlockedThreeSiteInjective, def:peps_edgeMiddleTensorFamily}
    If $A$ is vertex-injective and $e=(u,v)$ with $M_e=V\setminus\{u,v\}$, then
    $\inj(T_{A,M_e})\Longrightarrow
        \inj(T_{A,u})\wedge\inj(T_{A,M_e})\wedge\inj(T_{A,v})$.
\end{theorem}
\begin{proof}\leanok
    Combine $\inj(T_{A,M_e})$ with $\inj(T_{A,u})\wedge\inj(T_{A,v})$ from
    Definition~\ref{def:IsVertexInjective}.
\end{proof}

\begin{theorem}[All edge-blocked chains from middle-tensor injectivity]%
    \label{thm:peps_edgeBlockedThreeSiteInjective_all_of_middle}
    \lean{TNLean.PEPS.IsVertexInjective.edgeBlockedThreeSiteInjective_all_of_middle}
    \leanok
    \uses{def:IsVertexInjective, def:peps_edgeBlockedThreeSiteInjective, def:peps_edgeMiddleTensorFamily, thm:peps_edgeBlockedThreeSiteInjective_of_middle}
    If $A$ is vertex-injective and $\inj(T_{A,M_e})$ holds for every edge
    $e=(u,v)$, then
    $\forall e=(u,v),\ \inj(T_{A,u})\wedge\inj(T_{A,M_e})\wedge\inj(T_{A,v})$.
\end{theorem}
\begin{proof}\leanok
    Apply Theorem~\ref{thm:peps_edgeBlockedThreeSiteInjective_of_middle} at each
    edge $e=(u,v)$.
\end{proof}

\begin{theorem}[Region injectivity gives edge-blocked injectivity]%
    \label{thm:peps_edgeBlockedThreeSiteInjective_of_region}
    \lean{TNLean.PEPS.EdgeMiddleRegionInjectivityComparison.edgeBlockedThreeSiteInjective}
    \leanok
    \uses{def:IsVertexInjective, def:peps_edgeBlockedThreeSiteInjective, def:peps_edgeMiddleRegionInjectivityComparison, thm:peps_edgeBlockedThreeSiteInjective_of_middle}
    Let $e=(u,v)$ and $M_e=V\setminus\{u,v\}$. If $A$ is vertex-injective,
    $\kappa(M_e)$ holds, and $\kappa(M_e)\Longrightarrow\inj(T_{A,M_e})$, then
    $\inj(T_{A,u})\wedge\inj(T_{A,M_e})\wedge\inj(T_{A,v})$.
\end{theorem}
\begin{proof}\leanok
    Modus ponens gives $\inj(T_{A,M_e})$, and
    Theorem~\ref{thm:peps_edgeBlockedThreeSiteInjective_of_middle} concludes.
\end{proof}

\begin{theorem}[Middle tensors from all middle regions]%
    \label{thm:peps_edgeMiddleTensorInjective_all_of_region}
    \lean{TNLean.PEPS.EdgeMiddleRegionInjectivityComparison.edgeMiddleTensorInjective_all}
    \leanok
    \uses{def:peps_edgeMiddleRegionInjectivityComparison, def:peps_edgeMiddleTensorFamily}
    If $\kappa(M_e)$ and $\kappa(M_e)\Longrightarrow\inj(T_{A,M_e})$ hold for
    every edge $e=(u,v)$ with $M_e=V\setminus\{u,v\}$, then
    $\forall e,\ \inj(T_{A,M_e})$.
\end{theorem}
\begin{proof}\leanok
    Apply $\kappa(M_e)\Longrightarrow\inj(T_{A,M_e})$ to $\kappa(M_e)$ at each
    edge $e$.
\end{proof}

\begin{theorem}[All edge-blocked chains from middle-region injectivity]%
    \label{thm:peps_edgeBlockedThreeSiteInjective_all_of_region}
    \lean{TNLean.PEPS.EdgeMiddleRegionInjectivityComparison.edgeBlockedThreeSiteInjective_all}
    \leanok
    \uses{def:IsVertexInjective, def:peps_edgeBlockedThreeSiteInjective, def:peps_edgeMiddleRegionInjectivityComparison, thm:peps_edgeBlockedThreeSiteInjective_all_of_middle, thm:peps_edgeMiddleTensorInjective_all_of_region}
    If $A$ is vertex-injective and $\kappa(M_e)$ and
    $\kappa(M_e)\Longrightarrow\inj(T_{A,M_e})$ hold for every edge $e=(u,v)$,
    then $\forall e=(u,v),\ \inj(T_{A,u})\wedge\inj(T_{A,M_e})\wedge\inj(T_{A,v})$.
\end{theorem}
\begin{proof}\leanok
    Combine Theorem~\ref{thm:peps_edgeMiddleTensorInjective_all_of_region} with
    Theorem~\ref{thm:peps_edgeBlockedThreeSiteInjective_all_of_middle}.
\end{proof}

\begin{theorem}[Edge-middle tensor from singleton regions, nonempty middle case]%
    \label{thm:peps_edgeMiddleTensorInjective_of_singletons}
    \lean{TNLean.PEPS.EdgeMiddleRegionInjectivityComparison.edgeMiddleTensorInjective_of_singletons}
    \leanok
    \uses{thm:peps_finiteRegionInjectivityFromSingletons, def:peps_edgeMiddleRegionInjectivityComparison}
    Let $e=(u,v)$ and $M_e=V\setminus\{u,v\}\ne\varnothing$. If $\kappa$
    satisfies singleton comparison, union closure, and
    $\kappa(M_e)\Longrightarrow\inj(T_{A,M_e})$, and $A$ is vertex-injective,
    then $\inj(T_{A,M_e})$.
\end{theorem}
\begin{proof}\leanok
    Theorem~\ref{thm:peps_finiteRegionInjectivityFromSingletons} gives
    $\kappa(M_e)$, and $\kappa(M_e)\Longrightarrow\inj(T_{A,M_e})$ yields
    $\inj(T_{A,M_e})$.
\end{proof}

\begin{theorem}[One edge-blocked chain from singleton regions, nonempty middle case]%
    \label{thm:peps_edgeBlockedThreeSiteInjective_of_singletons}
    \lean{TNLean.PEPS.EdgeMiddleRegionInjectivityComparison.edgeBlockedThreeSiteInjective_of_singletons}
    \leanok
    \uses{thm:peps_edgeMiddleTensorInjective_of_singletons, thm:peps_edgeBlockedThreeSiteInjective_of_middle}
    Under the hypotheses of
    Theorem~\ref{thm:peps_edgeMiddleTensorInjective_of_singletons} at
    $e=(u,v)$, $\inj(T_{A,u})\wedge\inj(T_{A,M_e})\wedge\inj(T_{A,v})$.
\end{theorem}
\begin{proof}\leanok
    Combine Theorem~\ref{thm:peps_edgeMiddleTensorInjective_of_singletons} with
    Theorem~\ref{thm:peps_edgeBlockedThreeSiteInjective_of_middle}.
\end{proof}

\begin{theorem}[All edge-blocked chains from singleton regions, nonempty middle case]%
    \label{thm:peps_edgeBlockedThreeSiteInjective_all_of_singletons}
    \lean{TNLean.PEPS.EdgeMiddleRegionInjectivityComparison.edgeBlockedThreeSiteInjective_all_of_singletons}
    \leanok
    \uses{thm:peps_edgeBlockedThreeSiteInjective_of_singletons}
    If $M_e=V\setminus\{u,v\}\ne\varnothing$ and the hypotheses of
    Theorem~\ref{thm:peps_edgeMiddleTensorInjective_of_singletons} hold for
    every edge $e=(u,v)$, then
    $\forall e=(u,v),\ \inj(T_{A,u})\wedge\inj(T_{A,M_e})\wedge\inj(T_{A,v})$.
\end{theorem}
\begin{proof}\leanok
    Apply Theorem~\ref{thm:peps_edgeBlockedThreeSiteInjective_of_singletons} at
    each edge $e$.
\end{proof}

\begin{theorem}[All edge-blocked chains from singleton regions, cardinal form]%
    \label{thm:peps_edgeBlockedThreeSiteInjective_all_two_lt_card}
    \lean{TNLean.PEPS.EdgeMiddleRegionInjectivityComparison.edgeMiddleTensorInjective_all_two_lt_card}
    \lean{TNLean.PEPS.EdgeMiddleRegionInjectivityComparison.edgeBlockedThreeSiteInjective_all_two_lt_card}
    \leanok
    \uses{thm:peps_edgeMiddleVertices_nonempty_of_two_lt_card, thm:peps_edgeBlockedThreeSiteInjective_all_of_singletons, def:peps_edgeMiddleRegionInjectivityComparison}
    Assume singleton comparison, union closure for $\kappa$,
    $\forall e,\ \kappa(M_e)\Longrightarrow\inj(T_{A,M_e})$, and vertex
    injectivity of $A$. If $2<|V|$, then for every edge $e=(u,v)$,
    $\inj(T_{A,M_e})$ and
    $\inj(T_{A,u})\wedge\inj(T_{A,M_e})\wedge\inj(T_{A,v})$.
\end{theorem}
\begin{proof}\leanok
    For each $e=(u,v)$,
    Theorem~\ref{thm:peps_edgeMiddleVertices_nonempty_of_two_lt_card} gives
    $M_e=V\setminus\{u,v\}\ne\varnothing$, so
    Theorem~\ref{thm:peps_edgeBlockedThreeSiteInjective_all_of_singletons}
    concludes.
\end{proof}
